% Chain-parity note: a short standalone paper on lock parity, the mod-4 chain-count formula and the chain-parity law
% at a degree-5 vertex of a triangulated sphere. Draft, 8 Oct 2026. Separate from the VH-exists record
% (docs/reports/VHE-paper), on the coordinator's recommendation of 8 Oct (WK/CoordinatorPlan.md, "Paper decision").
% Label convention as in the VH-exists paper: every result carries one of [compiled] [hand] [computed] [cited] [open].
% Here [compiled] means: in Lean 4, no sorry, standard axioms only, and covered by audit J14 (WK/Audit-2026-10-08).
% Nothing is stated as a Theorem/Proposition/Corollary unless it is [compiled]; everything else carries its label.
% WK = SolvingFrameworkPlan/docs/working. Sources: WK/TrackF/LockParity.md; WK/TrackK/FProof.md (+TrackK-review);
% WK/TrackI/RigidIsolation.md (+TrackI-review); WK/TrackC/README.md sections 5-9; WK/Audit-2026-10-08/README.md;
% WK/TrackH-review/README.md (off-sphere rigid->rigid); WK/TrackS, WK/TrackT, WK/TrackU (matching form, n = 22 example).
% Open items for this draft: README.md beside this file.
\documentclass[11pt]{article}
\usepackage{amsmath,amssymb,amsthm,url,microtype}
\usepackage[hidelinks]{hyperref}
\usepackage[margin=1in]{geometry}
\usepackage{enumitem}\setlist{nosep}
\newtheorem{theorem}{Theorem}
\newtheorem{proposition}[theorem]{Proposition}
\newtheorem{corollary}[theorem]{Corollary}
\newtheorem{lemma}[theorem]{Lemma}
\theoremstyle{definition}
\newtheorem{definition}[theorem]{Definition}
\newtheorem*{example*}{Example}
\newcommand{\lab}[1]{\textsf{[#1]}}
\newcommand{\lean}[1]{\texttt{#1}}
\newcommand{\Z}{\mathbb Z}
\newcommand{\cw}{\mathrm{cw}}
\title{Kempe chains around a degree-five vertex:\\ a mod-4 chain count and a parity law}
% Working title was "A parity law for Kempe chains at a degree-five vertex". The proposed title names both results.
\author{Kyle Mathewson}
\date{8 October 2026 (draft)}
\sloppy
\begin{document}\maketitle

\begin{abstract}
This note, its proofs, its Lean formalisation and its computations were produced by AI agents (Claude models) under the direction of the author; Section~\ref{sec:disclosure} describes who did what and how it was checked.
Let $v$ be a vertex of degree five in a triangulation $T$ of the sphere, and $c$ a proper $4$-colouring of $T-v$ that uses all four colours around $v$. Let $N(c)$ be the number of Kempe chains of $c$, over all six pairs of colours. We give a formula for $N(c)$ mod~4 in terms of the number of faces of one Tait orientation, the order of $T$, the colours around $v$ and the two Kempe ``locks'' at $v$. Without $v$ it is a mod-4 form of Tutte's parity theorem combined with Fisk's degree. As a consequence, at a doubly locked colouring the exchange $\pi$ of Kempe's argument changes the parity of $N$ exactly when $\pi c$ is again doubly locked, so $\pi$ never maps a colouring with the minimum number of chains (eight) to another. We also show that each lock holds iff a certain Kempe chain contains an odd number of odd-degree vertices. All of this is formalised in Lean~4 and independently audited. Each result fails on the torus, and the law alone does not exclude near-rigid $\pi$-cycles. The Four Colour Theorem is known; this note gives no new proof of it.
\end{abstract}

\section{Introduction}\label{sec:intro}
Kempe's argument~\cite{kempe1879} colours $T-v$, for $v$ of degree at most five in a minimal counterexample, and frees a colour at $v$ by exchanging the colours of a bicoloured component, a \emph{Kempe chain}. Heawood~\cite{heawood1890} showed that at degree five the two exchanges can interfere. The Four Colour Theorem was proved by Appel and Haken~\cite{appel1977} and by Robertson, Sanders, Seymour and Thomas~\cite{robertson1997}, and checked formally by Gonthier~\cite{gonthier}. \textbf{This note claims no new proof of the Four Colour Theorem.} It studies the dynamics of Kempe's argument at a degree-five vertex, and proves parity constraints on them that hold on the sphere and fail on other surfaces.

\paragraph{Why these dynamics.} A colouring of $T-v$ is \emph{filled} if it extends to $v$. An unfilled colouring that is not doubly locked (Definition~\ref{def:locks}) becomes filled after one exchange. At a doubly locked colouring there is a canonical exchange $\pi$, and on the sphere $\pi$ permutes each Kempe class. So a Kempe class with no filled colouring is a union of $\pi$-cycles of doubly locked colourings. (Both facts are compiled in~\cite{vhe2026}, and are not used below.) Whether every Kempe class at every degree-five vertex has a filled colouring is, as far as we can tell, Tilley's open D-resolvability problem~\cite{tilley2017}. A companion record~\cite{vhe2026} reduces the Four Colour Theorem in Lean to a weaker, still open, statement of this kind. The results here (Section~\ref{sec:results}) constrain $\pi$-cycles on the sphere but do not exclude them. Section~\ref{sec:proofs} gives hand proofs, Section~\ref{sec:sharp} shows that each result fails off the sphere and that the law does not exclude near-rigid cycles, and Section~\ref{sec:lean} describes the Lean proofs and their audit.

\paragraph{Labels.} As in~\cite{vhe2026}, each result carries one label: \lab{compiled} (Lean~4, no \texttt{sorry}, standard axioms only, covered by audit J14); \lab{hand} (a written proof, reviewed by a different AI session unless marked otherwise); \lab{computed}; \lab{cited}; \lab{open}. Only \lab{compiled} results are called Theorem, Proposition or Corollary. No human has refereed any of this (Section~\ref{sec:disclosure}).

\section{Setting}\label{sec:setting}
A \emph{triangulation} $T$ is a simple graph embedded in the oriented sphere $S^2$ with every face bounded by a triangle; $n$ is its number of vertices. In Section~\ref{sec:proofs} we also allow parallel edges, but no loops, and ask only that every face be bounded by three distinct vertices. Each face is read as a cyclically ordered triple $(u,w,z)$ in the orientation of $S^2$.

\paragraph{The hole.} Let $v$ have degree five. Its neighbours form a $5$-cycle, labelled along the orientation so that the faces at $v$ are $(v,x_t,x_{t+1})$, indices mod~5. A \emph{state} is a proper colouring $c$ of $T-v$ with colours in $\Z_2^2=\{0,1,2,3\}$, with addition as bitwise exclusive or. It is \emph{filled} if the link uses at most three colours, and \emph{unfilled} otherwise. In an unfilled state exactly one colour repeats on the link, at two non-adjacent positions. After a cyclic shift of the indices, which preserves the orientation convention, the link reads
\[(c(x_0),\dots,c(x_4))=(\alpha,\beta,\alpha,\gamma,\delta),\qquad\{\alpha,\beta,\gamma,\delta\}=\Z_2^2,\]
the \emph{frame} of $c$. Unfilled states are always read in their own frame.

\paragraph{Kempe chains.} For colours $p\ne q$ the \emph{$pq$-chains} are the components of the subgraph of $T-v$ induced by the vertices coloured $p$ or $q$. $K_{pq}(x)$ is the vertex set of the chain containing $x$, and $p(p,q)$ is the number of $pq$-chains. A \emph{Kempe exchange} swaps $p$ and $q$ on one chain. The \emph{Kempe class} of $c$ is the set of states reachable from $c$ by exchanges. Let $N(c)=\sum_{\{p,q\}}p(p,q)$, summed over the six pairs, be the total number of Kempe chains.

\begin{definition}\label{def:locks}
Let $c$ be unfilled. Its \emph{locks} are $L_1$: $x_3\in K_{\beta\gamma}(x_1)$, and $L_2$: $x_4\in K_{\beta\delta}(x_1)$; we also write $L_i\in\{0,1\}$ for their indicators. The state is \emph{doubly locked} (DL) if $L_1=L_2=1$, and then $\pi c$ is obtained by exchanging $\alpha,\gamma$ on $K_{\alpha\gamma}(x_2)$. A DL state is \emph{rigid} if the pairs $\alpha\beta,\gamma\delta,\alpha\gamma,\beta\delta,\alpha\delta,\beta\gamma$ have $1,1,2,1,2,1$ chains.
\end{definition}
If $L_1$ fails, exchanging on $K_{\beta\gamma}(x_3)$ gives the filled link $(\alpha,\beta,\alpha,\beta,\delta)$. If $L_2$ fails, exchanging on $K_{\beta\delta}(x_4)$ gives $(\alpha,\beta,\alpha,\gamma,\beta)$. This is Kempe's step.

\paragraph{Orientation counts.} For a face $(u,w,z)$ the \emph{Tait triple} $(c(u){+}c(w),c(w){+}c(z),c(z){+}c(u))$ is an ordering of $\{1,2,3\}$. The face is \emph{clockwise} if the triple is a cyclic shift of $(1,2,3)$. Let $\cw(c)$ be the number of clockwise faces not containing $v$, and for unfilled $c$ let
\[\eta(c)=\big[(\alpha{+}\beta,\ \alpha{+}\gamma,\ \alpha{+}\delta)\text{ is a cyclic shift of }(1,2,3)\big]\in\{0,1\}.\]
For $X\subseteq V(T-v)$, $\delta(X)$ is the set of edges of $T$, edges to $v$ included, with exactly one end in $X$.

\section{Results}\label{sec:results}
Throughout this section $T$ is a triangulation of the sphere with $n$ vertices, $v$ is a vertex of degree five, and $c$ is an unfilled state. Lean names are listed in Section~\ref{sec:lean}.

\begin{proposition}[Kempe duality at the hole; \lab{compiled}]\label{prop:D}
$L_2$ holds iff $x_0\notin K_{\alpha\gamma}(x_2)$, and $L_1$ holds iff $x_0\notin K_{\alpha\delta}(x_2)$.
\end{proposition}

\begin{theorem}[Lock parity; \lab{compiled}]\label{thm:LP}
$L_2$ holds iff $|\delta(K_{\alpha\gamma}(x_2))|$ is odd, iff $K_{\alpha\gamma}(x_2)$ contains an odd number of vertices of odd degree in $T$. The same holds for $L_1$ with $K_{\alpha\delta}(x_2)$. Also, $|\delta(K_{\alpha\beta}(x_2))|$ is always odd.
\end{theorem}

\begin{proposition}[\lab{compiled}]\label{prop:eight}
If $c$ is doubly locked then $N(c)\ge 8$, with equality iff $c$ is rigid.
\end{proposition}

\begin{theorem}[Chain-count formula; \lab{compiled}]\label{thm:F}
$\displaystyle 2N(c)\equiv \cw(c)+(n-1)-\eta(c)+2\,(L_1+L_2)\pmod 4.$
\end{theorem}

\begin{theorem}[Chain-parity law; \lab{compiled}]\label{thm:law}
If $c$ is doubly locked, then $N(\pi c)+N(c)$ is odd iff $\pi c$ is doubly locked.
\end{theorem}

\begin{corollary}[Rigid isolation; \lab{compiled}]\label{cor:RI}
If $c$ is rigid, then $\pi c$ is not rigid.
\end{corollary}
Since $\pi c$ is one exchange away from $c$, it lies in the same Kempe class. Hence no Kempe class consists only of rigid states \lab{hand}; this one-line consequence is not stated separately in Lean. Along a $\pi$-cycle of DL states $N$ alternates in parity, so such cycles have even length \lab{hand}.

\begin{theorem}[\lab{compiled}]\label{thm:R7}
If $K$ is a chain of $c$ containing no neighbour of $v$, and $c'$ is obtained by the exchange on $K$, then $N(c')+L_1(c')+L_2(c')\equiv N(c)+L_1(c)+L_2(c)\pmod 2$.
\end{theorem}
$N$ alone is not invariant: on random spheres it changed parity in 111,912 of 281,230 such exchanges (\cite{record}, \texttt{TrackI}) \lab{computed}.

\begin{lemma}[Lemma W; \lab{compiled}]\label{lem:W}
If $c$ is doubly locked then $\cw(\pi c)\equiv\cw(c)+2\pmod 4$. An exchange on a chain that misses the link leaves $\cw$ unchanged mod~4; this part holds on every oriented triangulated surface.
\end{lemma}

\section{Proofs}\label{sec:proofs}
\subsection{Kempe duality and lock parity}
\begin{proof}[Proof of Proposition~\ref{prop:D}]
We treat $L_2$; for $L_1$ exchange $\gamma$ and $\delta$. If a $\beta\delta$-path $Q$ joins $x_1$ to $x_4$, then $C=vx_1Qx_4v$ is a cycle. The edges $vx_1,vx_4$ separate $vx_2,vx_3$ from $vx_0$ in the rotation at $v$, so $x_2$ and $x_0$ lie on different sides of $C$. An $\alpha\gamma$-path from $x_2$ to $x_0$ would miss $C$, contradicting the Jordan curve theorem.

Conversely, suppose $x_4\notin Q:=K_{\beta\delta}(x_1)$, and contract a spanning tree of the connected set $U=Q\cup\{v\}$ to one vertex $u$. In the resulting plane multigraph, two consecutive edges $ua,ub$ at $u$ bound a face coming from a face of $T$, so $a=b$ or $a\sim b$. Hence the neighbours of $u$ in rotation order form a closed walk. Near $v$ it reads $x_2,x_3,x_4,x_0$. Its other vertices are neighbours of $Q$, coloured $\alpha$ or $\gamma$, since a $\beta$- or $\delta$-coloured neighbour would lie in $Q$. In particular $x_4$ has no neighbour in $Q$ and occurs only once. The rest of the walk is an $\alpha\gamma$-walk in $T-v$ from $x_0$ to $x_2$.
\end{proof}

\begin{proof}[Proof of Theorem~\ref{thm:LP}]
By the handshake identity $\sum_{x\in X}\deg_T x=2e(T[X])+|\delta(X)|$, odd boundary and an odd number of odd-degree vertices are the same condition. By Proposition~\ref{prop:D} it then suffices to prove the following identity $(\ast)$, which uses no topology \lab{compiled for every triangulated rotation system}:
\[(\ast)\qquad|\delta(K_{\alpha\gamma}(x_2))|\text{ odd}\iff x_0\notin K_{\alpha\gamma}(x_2),\qquad|\delta(K_{\alpha\delta}(x_2))|\text{ odd}\iff x_0\notin K_{\alpha\delta}(x_2),\]
and $|\delta(K_{\alpha\beta}(x_2))|$ is odd.

Let $X=K_{pq}(x_2)$ and let $\{r,s\}$ be the complementary pair. (i) A face with one or two vertices in $X$ (\emph{mixed}) contains exactly two edges of $\delta(X)$, and every other face contains none. So $|\delta(X)|$ equals the number of mixed faces. (ii) Neighbours of $X$ other than $v$ are coloured $r$ or $s$. For $u\in X$ and $w\notin X\cup\{v\}$ adjacent, put $\tau(uw)=[c(u)=p]+[c(w)=r]\in\Z_2$. The two $\delta$-edges of a mixed face avoiding $v$ have different $\tau$: either the face's two vertices in $X$ are coloured $p,q$, or its two outside vertices are coloured $r,s$. Every $\delta$-edge not at $v$ lies in two faces avoiding $v$, except a link edge $x_tx_{t+1}$, which lies in one. Summing $\tau$ over mixed faces avoiding $v$ therefore gives
\[|\delta(X)|\equiv\#\{t:x_t\in X\text{ or }x_{t+1}\in X\}+\sum\nolimits_{t:\ |\{x_t,x_{t+1}\}\cap X|=1}\tau(x_tx_{t+1})\pmod 2 .\]
(iii) We evaluate this with $(p,q,r,s)=(\alpha,\gamma,\beta,\delta)$, $(\alpha,\delta,\beta,\gamma)$ and $(\alpha,\beta,\gamma,\delta)$; another labelling changes no parity. $K_{\alpha\gamma}(x_2)$ meets the link in $\{x_2,x_3\}$ or $\{x_0,x_2,x_3\}$, giving $3+0$ or $5+1$. $K_{\alpha\delta}(x_2)$ meets it in $\{x_2\}$ or $\{x_0,x_2,x_4\}$ (since $x_0x_4$ is an $\alpha\delta$-edge), giving $2+1$ or $5+1$. $K_{\alpha\beta}(x_2)$ meets it in $\{x_0,x_1,x_2\}$, giving $4+1$. Steps (i)--(ii) use only that every edge lies in two faces and that the faces at $v$ are the $(v,x_t,x_{t+1})$.
\end{proof}

\begin{proof}[Proof of Proposition~\ref{prop:eight}]
At a DL state, Proposition~\ref{prop:D} gives $x_0\notin K_{\alpha\gamma}(x_2)$ and $x_0\notin K_{\alpha\delta}(x_2)$, so these pairs have at least two chains each. Each other pair has a colour on the link, so it has at least one chain. Hence $N\ge 8$, with equality exactly for the rigid counts.
\end{proof}

\subsection{Tutte's identity and Fisk's degree}
Here $T$ is a loopless triangulation of the sphere, parallel edges allowed, with $F=2n-4$ faces, and $c$ is a proper $4$-colouring of all of $T$. Let $|X|$ be the number of vertices coloured $X$, and $e(X,Y)$ the number of $XY$-edges.

\begin{lemma}[Tutte~\cite{tutte1969,mohar2019}; \lab{cited}; compiled in a sides form; proof written for this note]\label{lem:tutte}
If $\{X,Y,Z,W\}$ are the four colours, then $p(X,Y)-p(Z,W)=|X|+|Y|-e(X,Y)-1$.
\end{lemma}
\begin{proof}
By Euler's formula, the plane multigraph $G=T_{XY}$, which has $p(X,Y)$ components, has $1+p(X,Y)-|X|-|Y|+e(X,Y)$ faces. Each $ZW$-chain is connected and misses $G$, so it lies in one face of $G$. Each face $R$ of $G$ contains a face of $T$, whose corner coloured $Z$ or $W$ lies in $R$. The $ZW$-corners of one $T$-face are adjacent. Two $T$-faces in $R$ that share an edge outside $G$ share a $ZW$-coloured end of that edge, and $R$ minus its vertices is connected through such edges. So each face of $G$ contains exactly one $ZW$-chain.
\end{proof}

\begin{lemma}[Fisk's degree~\cite{fisk1977,mohar2019}, mod 4; \lab{hand}, independent AI review]\label{lem:fisk}
On a closed oriented loopless triangulated surface with $F$ faces, every proper $4$-colouring satisfies $\cw\equiv F/2+2\deg(A)\pmod 4$ for each colour $A$, where $\deg(A)=\sum_{c(a)=A}\deg a$.
\end{lemma}
\begin{proof}
Let $\sigma(f)=\pm1$ according as $f$ is clockwise or not. If $f$ reads colours $(a,b,e)$ and $d$ is the fourth colour, then $\sigma(f)=+1$ iff $(d,a,b,e)$ is an even permutation of $(0,1,2,3)$; this is a check of 24 cases. So two faces reading an edge as $(a,b,z)$ and $(b,a,w)$ have equal $\sigma$ iff $z\ne w$. Let $d_t=\sum_{f\text{ coloured }t}\sigma(f)$ for each colour triangle $t$. For $t=\{x,y,z\}$ and $t'=\{x,y,w\}$, pairing faces across their $xy$-edges gives $d_t=d_{t'}$: two faces with the same third colour cancel, and faces with third colours $z,w$ contribute equally to both. So all $d_t$ equal one integer $d$, and $2\cw-F=\sum_t d_t=4d$. Every face has at most one corner coloured $A$, so $F-\deg(A)$ faces are coloured by the triangle avoiding $A$, and $d\equiv F-\deg(A)\equiv\deg(A)\pmod 2$.
\end{proof}

\begin{lemma}[\lab{hand}; the formula with independent AI review, the last sentence new here and unreviewed]\label{lem:F0}
$2N\equiv \cw+n\pmod 4$. A Kempe exchange leaves $\cw$ unchanged mod~$4$, on every closed oriented triangulated surface.
\end{lemma}
\begin{proof}
Fix $A$ and sum Lemma~\ref{lem:tutte} over the three pairs $\{A,X\}$. Since $N\equiv\sum(p(A,X)-p(\text{complement}))\pmod 2$ and $\sum_X e(A,X)=\deg(A)$, this gives $N\equiv 3|A|+(n-|A|)+\deg(A)+3\equiv n+1+\deg(A)\pmod 2$, which is Tutte's parity theorem~\cite[Thms.~1, 3]{mohar2019}. With Lemma~\ref{lem:fisk}, $\cw+n\equiv 2n-2+2\deg(A)\equiv 2N\pmod 4$. An exchange of $p,q$ fixes the vertices of each colour $A\notin\{p,q\}$, hence $\deg(A)$ and $\cw$ mod~$4$.
\end{proof}

\subsection{The hole: filling the pentagon}
\begin{proof}[Proof of Theorem~\ref{thm:F}]
Delete $v$, and add the diagonals $x_1x_3$ and $x_1x_4$ inside the resulting pentagonal face, as new edges even if $T$ already has such an edge. The result $T^\circ$ is a loopless triangulation of the sphere with $n-1$ vertices and new faces $(x_1,x_2,x_3)$, $(x_1,x_3,x_4)$, $(x_4,x_0,x_1)$. The new edges are coloured $\beta\gamma$ and $\beta\delta$, so $c$ is proper on $T^\circ$.
\emph{Chains:} an added edge merges two chains iff its ends were in different chains, and the diagonals are exactly the lock edges, so $N(c)=N_{T^\circ}(c)+(1-L_1)+(1-L_2)$.
\emph{Faces:} the new faces read $(\beta,\alpha,\gamma)$, $(\beta,\gamma,\delta)$, $(\delta,\alpha,\beta)$. With $a=\alpha{+}\beta$, $g=\alpha{+}\gamma$, $e=\alpha{+}\delta$, their Tait triples are $(a,g,e)$, $(e,a,g)$ and $(e,a,g)$, each clockwise iff $\eta=1$. So $\cw_{T^\circ}(c)=\cw(c)+3\eta$.
Lemma~\ref{lem:F0} for $T^\circ$ gives $2N(c)\equiv\cw(c)+3\eta+(n-1)+4-2L_1-2L_2\equiv\cw(c)+(n-1)-\eta+2(L_1+L_2)\pmod 4$.
\end{proof}
Labelling the link against the orientation would turn $\eta$ into $1-\eta$ and break the formula; Lean reads $\eta$ along the face orientation (\lean{handS}).

\begin{proof}[Proof of Lemma~\ref{lem:W} \lab{hand, new in this note and unreviewed; Lean uses a dart sum instead}]
Let $c$ be DL and $K=K_{\alpha\gamma}(x_2)$. Then $x_3\in K$, $x_1,x_4\notin K$, and $x_0\notin K$ by Proposition~\ref{prop:D}, so $c'=\pi c$ has link $(\alpha,\beta,\gamma,\alpha,\delta)$. The diagonals of $T^\circ$ are not $\alpha\gamma$-edges, so $K$ is still a chain of $T^\circ$, and by Lemma~\ref{lem:F0}, $\cw_{T^\circ}(c')\equiv\cw_{T^\circ}(c)$. Under $c'$ the fill faces read $(\beta,\gamma,\alpha)$, $(\beta,\alpha,\delta)$, $(\delta,\alpha,\beta)$, with Tait triples $(e,g,a)$, $(a,e,g)$, $(e,a,g)$, clockwise iff $\eta=0,0,1$. Hence $\cw(c')+2-\eta\equiv\cw(c)+3\eta$, so $\cw(c')\equiv\cw(c)+2\pmod 4$. For a chain missing the link the fill faces do not change.
\end{proof}

\begin{proof}[Proof of Theorem~\ref{thm:law}]
Let $c$ be DL and $c'=\pi c$. Read from $x_3$, the link of $c'$ is $(\alpha,\delta,\alpha,\beta,\gamma)$, so the roles in its frame are $(\alpha',\beta',\gamma',\delta')=(\alpha,\delta,\beta,\gamma)$. Then $\eta(c')=\eta(c)$, since $(e,a,g)$ is a cyclic shift of $(a,g,e)$. Also $L_1(c')=[x_1\in K_{\beta\delta}(x_4)]=L_2(c)=1$, because an $\alpha\gamma$-exchange does not change the $\beta\delta$-chains. So $c'$ is DL iff $L_2(c')=1$. Subtracting Theorem~\ref{thm:F} at $c$ from Theorem~\ref{thm:F} at $c'$ and using Lemma~\ref{lem:W} gives $2(N(c')-N(c))\equiv 2+2(1+L_2(c'))-4\equiv 2L_2(c')\pmod 4$.
\end{proof}

\begin{proof}[Proof of Corollary~\ref{cor:RI} and Theorem~\ref{thm:R7}]
If $c$ and $\pi c$ were both rigid, both would have $N=8$ by Proposition~\ref{prop:eight}, while Theorem~\ref{thm:law} requires an odd sum. For Theorem~\ref{thm:R7}: the exchange fixes the link colours, hence $\eta$, and fixes $\cw$ mod~4 (Lemma~\ref{lem:W}), so Theorem~\ref{thm:F} gives $2\Delta N\equiv2\Delta(L_1+L_2)\pmod 4$.
\end{proof}

\subsection{A second proof of the law by curve surgery (sketch)}\label{sec:surgery}
\lab{hand}, independent AI review; not formalised. In the dual of $T$, contract the five faces at $v$ to a vertex $v$ of degree five; the other vertices are cubic. Colour the edge dual to $uw$ by $c(u)+c(w)$, with $1=\alpha{+}\beta$, $2=\alpha{+}\gamma$, $3=\alpha{+}\delta$. The edges $e_t$ at $v$, dual to $x_tx_{t+1}$, are then coloured $1,1,2,1,3$. Let $M_i$ be colour class $i$, and put $H=M_2\cup M_3$, $F_{12}=M_1\cup M_2$, $F_{13}=M_1\cup M_3$. By Euler's formula, the regions of the complements of these three subgraphs are the chains of the pairs $\{\alpha\beta,\gamma\delta\}$, $\{\alpha\delta,\beta\gamma\}$ and $\{\alpha\gamma,\beta\delta\}$ respectively, so $N=5+k(H)+k(F_{12})+k(F_{13})$, where $k$ counts components. By the Jordan curve theorem, $L_1$ and $L_2$ are the non-crossing pairings $(e_0e_3)(e_1e_2)$ of $F_{12}$ and $(e_1e_3)(e_0e_4)$ of $F_{13}$ at $v$. The exchange $\pi$ swaps colours $1$ and $3$ on the boundary $C$ of the region of $K_{\alpha\gamma}(x_2)$. This boundary consists of the $F_{13}$-loop through $e_1,e_3$ and some $F_{13}$-cycles, and every other complementary region of $C$ is a disc. In the frame of $\pi c$ one has $H'=F_{12}\triangle C$, $F_{12}'=F_{13}$ and $F_{13}'=H\triangle C$. The law becomes $k(H)+k(F_{12})+k(H\triangle C)+k(F_{12}\triangle C)\equiv[H\triangle C\text{ pairs }(e_1e_4)(e_2e_3)]\pmod 2$. This is a band-surgery parity lemma: non-crossing matchings in a disc are connected by flips, each flip acts on both curve systems, and on the sphere a band move changes the number of curves by exactly $\pm1$ (Schoenflies).

\section{Sharpness}\label{sec:sharp}
\paragraph{Off the sphere.} Each result fails on other surfaces, and only at the steps that use the sphere.
\emph{Lock parity:} $(\ast)$ held at all 21,968,170 unfilled states tested on random triangulations of the torus, the Klein bottle, the projective plane, and the surfaces of orientable genus~2 and non-orientable genus~3 \lab{computed}; Proposition~\ref{prop:D} fails, e.g.\ at 1,664,226 of 3,862,132 unfilled torus states \lab{computed}.
\emph{Chain-count formula:} for each partition of the colours into two pairs, the dual edges of the edges of $T^\circ$ that join the two classes form a Tait $2$-factor. On the sphere its complementary regions are planar, one per chain, which is the region form of Lemma~\ref{lem:tutte}. On an orientable surface, let $G$ be the total genus of these regions over the three partitions. The argument then gives
\[2N-\cw-(n-1)+\eta-2(L_1+L_2)\equiv 2G\pmod 4\]
\lab{hand; this genus extension is not independently reviewed}. On the torus the residue was $2G$ at all 2,192 hole states tested, with $G=0,1,2$ at $340$, $1{,}137$ and $715$ of them. So the formula fails at the 1,137 states with $G$ odd \lab{computed}. In Lean the sphere enters the formula only through \lean{Fills}, which fails on the torus.
\emph{Law and rigid isolation:} on the torus the law failed at 282 of 544 $\pi$-steps \lab{computed}. Rigid states mapped to rigid states 23 times among 2,467 rigid states on the torus, once among 1,401 on the Klein bottle and 29 times among 5,205 on the projective plane; a second program confirmed all 53 \lab{computed}.

\paragraph{The law does not exclude near-rigid loops.} Call a $\pi$-cycle of DL states \emph{near-rigid} if $N\le 9$ throughout. By Proposition~\ref{prop:eight} and the law, $N$ then alternates $8,9,8,9,\dots$; rigid isolation excludes two consecutive 8s. In the dual picture of Section~\ref{sec:surgery}, along DL states with $F_{12}$, $F_{13}$ connected, a state is determined by the perfect matching $M_t=M_2$. The next matching is the previous one switched along the odd loop of the figure-eight $E-M_t$, and $N=7+k(M_{t-1}\cup M_t)$ \lab{hand, not independently reviewed}. These conditions make sense in any graph that is cubic except for one degree-five vertex with a cyclic order on its edges. On the sphere every near-rigid cycle yields such a sequence with $k(M_{t-1}\cup M_t)$ alternating $1,2$ \lab{hand, not independently reviewed}. Relabelling the edges at $v$ of plane examples in each of the 22 non-planar cyclic orders makes the law fail at 611 of 1,261 steps, against 0 of 256 for the planar order \lab{computed}.

\begin{example*}[\lab{computed}, two independent programs]
A $3$-connected non-planar graph on $22$ vertices, cubic except for one vertex of degree five, carries a cyclic sequence of 10 perfect matchings that satisfies all of these conditions. Along it $k(M_{t-1}\cup M_t)=1,2,1,2,\dots$ and the abstract chain-parity law holds at every step. It was found by shrinking an order-28 example; none was found at order 20, and none is known to be dual to a surface triangulation. The edge list is in~\cite{record}, \texttt{TrackU/README.md}.
\end{example*}
So the parity law, together with the local matching structure, does not exclude near-rigid cycles. A proof must use planarity again.

\paragraph{Open problems \lab{open}.} \emph{NRC}: on a triangulated sphere, no $\pi$-cycle at a degree-five vertex consists of DL states with $N\le9$. \emph{LPC}: no Kempe class of $T-v$ consists only of states on DL $\pi$-cycles; at every degree-five vertex this is, as far as we can tell, Tilley's D-resolvability~\cite{tilley2017}. In an exhaustive census of a restricted class of triangulations of orders 22--33, no near-rigid cycle occurs, but the longest near-rigid run along a $\pi$-orbit grows with the order and reaches 8 at order 33~\cite{vhe2026} \lab{computed}.

\section{Formal verification}\label{sec:lean}
The Lean~4 library~\cite{planemap} models a map as a \lean{SphericalMap}: a simple graph on \lean{Fin n}, a rotation system, and \lean{Fills}, which says that every even edge set is a $\Z_2$-sum of face boundaries. \lean{Triangulated} means every face has length~3. Connectivity is not assumed, and the bridge from embedded graphs to this model is \emph{not} formalised. The hole is \lean{Pent}, a frame is \lean{RepeatAt}, the locks are \lean{Lock1} and \lean{Lock2}, DL states are \lean{DoublyLocked} (or \lean{DLState}, DL in some frame), and $\pi$ is \lean{piMove}. The formula and its corollaries assume no isolated vertex. Namespace: \lean{SimpleGraph.QuarterFloor} (Tutte: \lean{SimpleGraph.TutteSides}).

\begin{center}\small
\begin{tabular}{@{}p{2.3cm}p{13.4cm}@{}}
Result & Lean names (module)\\\hline
Prop.~\ref{prop:D} & \lean{lock2\_iff\_not\_reach\_alphaA}, \lean{lock1\_iff\_not\_reach\_alphaB} (\lean{QuarterLockParity})\\
Thm.~\ref{thm:LP} & \lean{lock2\_iff\_odd\_boundary}, \lean{lock1\_iff\_odd\_boundary}, \lean{alphaMu\_odd\_boundary} and the three \lean{\_oddCount} forms; $(\ast)$: \lean{parity\_alphaA}, \lean{parity\_alphaB}, \lean{parity\_alphaMu} (\lean{QuarterLockParity})\\
Prop.~\ref{prop:eight} & \lean{eight\_le\_nChains}, \lean{rigidAt\_iff}; definitions \lean{chainCount}, \lean{nChains}, \lean{RigidAt} (\lean{ChainCount})\\
Lemma~\ref{lem:W} & \lean{cw\_piMove} (sphere); \lean{lemmaW}, \lean{lemmaW\_linkFree}, \lean{cw\_swap\_mod8} (any triangulated rotation system) (\lean{ChainMod4})\\
Lemma~\ref{lem:tutte} & \lean{tutte\_sides}, \lean{euler\_tri} (\lean{TutteSides})\\
Thm.~\ref{thm:F} & \lean{chainFormulaF}, \lean{conjectureF} (\lean{ChainF}); statements \lean{ChainFormulaF}, \lean{ConjectureF} (\lean{ChainMod4})\\
Thm.~\ref{thm:law} & \lean{chainParityLaw\_sphere} (\lean{ChainF}); statement \lean{ChainParityLaw} (\lean{ChainCount}); \lean{chainParityLaw\_of\_F} (\lean{ChainMod4})\\
Cor.~\ref{cor:RI}, Thm.~\ref{thm:R7} & \lean{rigid\_isolation}, \lean{remark7} (\lean{ChainF})\\
\end{tabular}
\end{center}

\paragraph{The Lean route.} Tutte's identity is proved by rank computations over $\Z_2$; \lean{face\_side} is the only use of \lean{Fills}. Instead of filling the pentagon, $v$ is put on one side of each colour partition, its two edges there playing the lock diagonals, and $\cw$ mod~4 is a dart sum against tables checked by \lean{decide} (\lean{omT\_face}, \lean{omT\_link}; no \texttt{native\_decide}).

\paragraph{Audit.} Audit J14~\cite{audit}, of 8 October 2026, was run by a session that wrote none of the code. It worked at commit \texttt{ad0fd00b}, with Mathlib \texttt{300d0e5} and Lean \texttt{v4.35.0-rc3}. It rebuilt all 91 modules from source (exit 0, no \texttt{sorry}), and checked that every \texttt{\#print axioms} line lies within \texttt{propext}, \texttt{Classical.choice}, \texttt{Quot.sound}; a planted \texttt{sorry} was caught. It compared each statement above with its prose and instantiated the hypotheses of Theorems~\ref{thm:LP}, \ref{thm:F}, \ref{thm:law}, Proposition~\ref{prop:eight} and Lemma~\ref{lem:W} at an explicit DL state of a $22$-vertex triangulation (\texttt{NonVacuity.lean}). (There $N(c)=N(\pi c)=12$ and $\pi c$ is not DL. Rigid states were not exhibited.) The audit's wording notes are applied here: Lemma~\ref{lem:W} at $\pi$ is sphere-only, the class form of rigid isolation is \lab{hand}, and Theorem~\ref{thm:R7} is stated at an unfilled state. The Lean authors also checked that sabotaged variants (e.g.\ $n$ for $n-1$, or the sign of $\eta$ flipped) fail to compile. \emph{Frozen statements:} unlike the reductions in~\cite{vhe2026}, these results have no frozen challenge file with a proved bridge yet. The audited file hashes (SHA-256) begin \texttt{956a1956} (\lean{QuarterLockParity}), \texttt{ba3184b2} (\lean{ChainCount}), \texttt{8c9cc088} (\lean{ChainMod4}), \texttt{278a9ee9} (\lean{TutteSides}) and \texttt{6d87b43d} (\lean{ChainF}).

\paragraph{Reproducing.} From the record~\cite{record}, with a built PlaneMap checkout at Mathlib \texttt{300d0e5}:
\begin{verbatim}
cd SolvingFrameworkPlan/docs/working/StudioMathLean
shasum -a 256 -c SHA256SUMS
./check.sh <built-planemap-checkout> <fresh-scratch-dir>
\end{verbatim}
\texttt{check.sh} compiles all modules in dependency order and prints the axioms of each headline theorem. The audit's scripts and \texttt{NonVacuity.lean} are in \texttt{docs/working/Audit-2026-10-08}.

% ===== Disclosure: adapted from the VH-exists paper (8-9 Oct); Kyle approved the substance (9 Oct), final text to be shown to him before posting.
% KYLE MUST APPROVE THIS TEXT FOR THIS NOTE before it is posted. It was written for the VH-exists paper and still says "this paper"
% and refers to the Navigator's ledger, AUTHOR-RECORD.md "beside this paper", and session roles; adapt only with Kyle's approval.
\section{How this work was produced}\label{sec:disclosure}
{\small
\textbf{The mathematics, the Lean proofs, the computations and the text of this note were produced by AI agents, not by a human.} The author directed the agents; he did not write the proofs, the code or this text, and he has not checked them line by line. Throughout the note, ``we'' means these AI agents working under the author's direction. The agents were Claude models made by Anthropic, run in Claude Code: long-running sessions that read and wrote files, ran programs on the author's computers, and exchanged dated messages through a shared repository. The repository's commit trailers name the models used. Every commit carries Kyle Mathewson as its git author because the agents committed under his identity; the trailers, not the author field, show who wrote each change.

\textbf{How these results were produced.} The work behind this note was done on 7--8 October 2026, inside a larger programme described in~\cite{vhe2026}. On 7 October the author made one Claude Code session the sole coordinator of that programme and delegated to it the choice of tasks and computing budget. The coordinator ran short-lived sub-agents with fixed, written briefs: some searched for counterexamples, some wrote hand proofs, some wrote the Lean~4 formalisation, and others, which had not seen the work under review, refereed hand proofs and audited the Lean modules with their own independently written code. The coordinator reran the Lean regression itself before accepting any module. The earlier team sessions of the programme (Long Table, Math, Audit, Navigator, and a night swarm of exploratory agents) are described in~\cite{vhe2026}; some lemmas used here first appeared in their notes.

\textbf{What the author did.} Kyle Mathewson is responsible for the claims. He set the goal of a new, deductive route to the Four Colour Theorem, appointed the coordinator and delegated authority to it, provided the machines (a MacBook Pro and a Mac Studio) on which every computation ran, approved installing outside software, asked for the visual and expository material, and decided that this result be written up as a separate note, its title, its authorship and this disclosure. \texttt{AUTHOR-RECORD.md} beside~\cite{vhe2026} lists his actions with their sources. He did not review the hand proofs, the Lean statements or the computed outputs.

\textbf{How the results were checked, and by whom.} All checking was done by AI sessions or by the Lean kernel; none was done by a human. Every claim carries one label.
\begin{itemize}
\item \lab{compiled}: proved in Lean~4 with no \texttt{sorry} and only Lean's standard axioms, and covered by the independent module audit J14~\cite{audit} (Section~\ref{sec:lean}). The Lean kernel certifies that the formal statement is proved. It does not certify that the formal statement means what the prose in this note says; that match was checked only by AI sessions (the audit compared each statement with the text and showed the hypotheses are satisfiable), and it is the main remaining risk in trusting a \lab{compiled} result.
\item \lab{hand}: a written proof reviewed by a different AI session. No human referee has checked any of these. AI reviewers can share blind spots, so treat these as unrefereed.
\item \lab{computed}: a computation whose output was checked by a second, separately written implementation, and which claims only what happened on the graphs it ran.
\item \lab{cited}: taken from the published literature as summarised by the agents; consult the original.
\item \lab{open}: not claimed.
\end{itemize}
Mistakes were made along the way, among them a lemma first stated in a false form and conjectures later refuted. Each was corrected in a dated record rather than erased.

\textbf{The record.} The full record is public at \url{https://github.com/kylemath/GraphColourDeductiveSolver}, and the Lean development at \url{https://github.com/kylemath/mathlib4-planemap}. A tagged release, [\textsc{tag name to be set at posting}], fixes the version this note describes. The repository holds the coordinator's plan and status log, every review and audit report, every declaration with its hash, and all code and outputs. It also holds fiction written by the agents, which is not part of the mathematical record.

}
\begin{footnotesize}
\begin{thebibliography}{99}\setlength{\itemsep}{0pt}
\bibitem{appel1977} K.~Appel and W.~Haken, Every planar map is four colorable. Part~I: Discharging, \emph{Illinois J. Math.} 21(3) (1977) 429--490.
\bibitem{audit} Audit J14: the 7--8 October Lean modules and the frozen challenges, \texttt{SolvingFrameworkPlan/docs/working/Audit-2026-10-08/README.md} in~\cite{record} (8 October 2026).
\bibitem{fisk1977} S.~Fisk, Geometric coloring theory, \emph{Adv. Math.} 24 (1977) 298--340.
\bibitem{gonthier} G.~Gonthier, Formal proof---the four-color theorem, \emph{Notices Amer. Math. Soc.} 55(11) (2008) 1382--1393.
\bibitem{heawood1890} P.~J.~Heawood, Map-colour theorem, \emph{Quart. J. Pure Appl. Math.} 24 (1890) 332--338.
\bibitem{kempe1879} A.~B.~Kempe, On the geographical problem of the four colours, \emph{Amer. J. Math.} 2(3) (1879) 193--200.
\bibitem{mohar2019} B.~Mohar and N.~Singer, The last temptation of William T.~Tutte, arXiv:1912.07205 (2019); to appear in \emph{European J. Combin.} (per arXiv).
\bibitem{planemap} The PlaneMap library for Lean~4 and Mathlib, \url{https://github.com/kylemath/mathlib4-planemap} (snapshot \texttt{8299419}).
\bibitem{record} The public record (repository): \url{https://github.com/kylemath/GraphColourDeductiveSolver}.
\bibitem{robertson1997} N.~Robertson, D.~P.~Sanders, P.~Seymour and R.~Thomas, The four-colour theorem, \emph{J. Combin. Theory Ser. B} 70(1) (1997) 2--44.
\bibitem{tilley2017} J.~Tilley, D-resolvability of vertices in planar graphs, \emph{J. Graph Algorithms Appl.} 21(4) (2017) 649--661, doi:10.7155/jgaa.00433.
\bibitem{tutte1969} W.~T.~Tutte, Even and odd 4-colorings, in \emph{Proof Techniques in Graph Theory}, Academic Press, 1969, 161--169.
\bibitem{vhe2026} K.~Mathewson, The vacancy induction for the Four Colour Theorem: a dated record of what is solid, draft (October 2026), \texttt{SolvingFrameworkPlan/docs/reports/VHE-paper} in~\cite{record}.
\end{thebibliography}
\end{footnotesize}
\end{document}
